\documentclass[11pt,a4paper]{article}
\usepackage[utf8]{inputenc}
\usepackage[margin=1in,headheight=14pt]{geometry}
\usepackage{hyperref}
\usepackage{listings}
\usepackage{xcolor}
\usepackage{graphicx}
\usepackage{fancyhdr}
\usepackage{tcolorbox}
\usepackage{amssymb}

% Code listing style
\lstset{
    basicstyle=\ttfamily\small,
    breaklines=true,
    frame=single,
    backgroundcolor=\color{gray!10},
    keywordstyle=\color{blue},
    commentstyle=\color{green!60!black},
    stringstyle=\color{red},
    showstringspaces=false
}

% Header and footer
\pagestyle{fancy}
\fancyhf{}
\rhead{Drone Detection Model Comparison}
\lhead{Deployment Guide}
\rfoot{Page \thepage}

\title{\textbf{Drone Detection Model Comparison App\\Deployment Guide}}
\author{VisionSentry Project}
\date{\today}

\begin{document}

\maketitle
\tableofcontents
\newpage

\section{Introduction}

This guide provides step-by-step instructions for deploying the Drone Detection Model Comparison web application. The application uses YOLOv11/12 models for drone detection and provides a web interface for comparing two models side-by-side.

\subsection{Key Features}
\begin{itemize}
    \item Side-by-side model comparison
    \item Synchronized video playback
    \item Real-time detection visualization
    \item Detection statistics and graphs
    \item CPU-optimized for cloud deployment
\end{itemize}

\subsection{System Requirements}
\begin{itemize}
    \item Python 3.11+
    \item 2GB+ RAM (4GB recommended)
    \item CPU-only (no GPU required)
    \item Docker (for containerized deployment)
\end{itemize}

\section{Deployment Options}

\subsection{Why Not Vercel?}

\begin{tcolorbox}[colback=red!5!white,colframe=red!75!black,title=Important Note]
\textbf{Vercel does not support Python applications.} Vercel is designed for JavaScript/Node.js frameworks (Next.js, React, etc.) and static sites only. For Python ML applications, use the alternatives below.
\end{tcolorbox}

\subsection{Recommended Platforms}

\begin{enumerate}
    \item \textbf{Hugging Face Spaces} - Free, ML-optimized, easiest deployment
    \item \textbf{Railway.app} - Simple, free tier available
    \item \textbf{Render.com} - Free tier, Docker support
    \item \textbf{Google Cloud Run} - Pay-per-use, scalable
    \item \textbf{AWS EC2} - Full control, requires more setup
\end{enumerate}

\section{Option 1: Hugging Face Spaces (Recommended)}

\subsection{Why Hugging Face Spaces?}
\begin{itemize}
    \item \textbf{Free tier} with generous resources
    \item Optimized for ML applications
    \item Built-in Gradio support
    \item No credit card required
    \item Automatic HTTPS
\end{itemize}

\subsection{Deployment Steps}

\subsubsection{Step 1: Create Account}
\begin{enumerate}
    \item Go to \url{https://huggingface.co}
    \item Sign up for a free account
    \item Verify your email
\end{enumerate}

\subsubsection{Step 2: Create New Space}
\begin{enumerate}
    \item Click on your profile → ``Spaces''
    \item Click ``Create new Space''
    \item Configure:
    \begin{itemize}
        \item Space name: \texttt{drone-detection-comparison}
        \item License: MIT or Apache 2.0
        \item SDK: \textbf{Gradio}
        \item Hardware: CPU basic (free)
    \end{itemize}
    \item Click ``Create Space''
\end{enumerate}

\subsubsection{Step 3: Upload Files}
Upload the following files to your Space:
\begin{lstlisting}[language=bash]
model_comparison_app.py
requirements_app.txt
weights/best.pt
VisionSentry-dex-rgb-model-replication/weights/best.pt
README.md
\end{lstlisting}

\subsubsection{Step 4: Configure App File}
Rename \texttt{model\_comparison\_app.py} to \texttt{app.py} (Hugging Face convention):
\begin{lstlisting}[language=bash]
mv model_comparison_app.py app.py
\end{lstlisting}

\subsubsection{Step 5: Deploy}
\begin{enumerate}
    \item Commit files to the Space repository
    \item Hugging Face will automatically build and deploy
    \item Wait 2-5 minutes for deployment
    \item Access your app at: \texttt{https://huggingface.co/spaces/YOUR\_USERNAME/drone-detection-comparison}
\end{enumerate}

\section{Option 2: Railway.app}

\subsection{Deployment Steps}

\subsubsection{Step 1: Setup}
\begin{enumerate}
    \item Go to \url{https://railway.app}
    \item Sign up with GitHub
    \item Click ``New Project'' → ``Deploy from GitHub repo''
\end{enumerate}

\subsubsection{Step 2: Configure}
\begin{enumerate}
    \item Select your repository
    \item Railway will detect the Dockerfile
    \item Set environment variables (if needed)
    \item Click ``Deploy''
\end{enumerate}

\subsubsection{Step 3: Access}
\begin{enumerate}
    \item Wait for build to complete (5-10 minutes)
    \item Railway will provide a public URL
    \item Access your app at the provided URL
\end{enumerate}

\section{Option 3: Render.com}

\subsection{Deployment Steps}

\subsubsection{Step 1: Create Account}
\begin{enumerate}
    \item Go to \url{https://render.com}
    \item Sign up with GitHub
\end{enumerate}

\subsubsection{Step 2: Create Web Service}
\begin{enumerate}
    \item Click ``New +'' → ``Web Service''
    \item Connect your GitHub repository
    \item Configure:
    \begin{itemize}
        \item Name: \texttt{drone-detection}
        \item Environment: Docker
        \item Plan: Free
    \end{itemize}
\end{enumerate}

\subsubsection{Step 3: Deploy}
\begin{enumerate}
    \item Render will use the \texttt{render.yaml} configuration
    \item Click ``Create Web Service''
    \item Wait for deployment (10-15 minutes first time)
\end{enumerate}

\section{Local Deployment}

\subsection{Using Docker}

\subsubsection{Build Image}
\begin{lstlisting}[language=bash]
cd VisionSentry
docker build -t drone-detection-app .
\end{lstlisting}

\subsubsection{Run Container}
\begin{lstlisting}[language=bash]
docker run -p 7860:7860 drone-detection-app
\end{lstlisting}

\subsubsection{Access}
Open browser to: \texttt{http://localhost:7860}

\subsection{Using Python Directly}

\subsubsection{Install Dependencies}
\begin{lstlisting}[language=bash]
pip install -r requirements_app.txt
\end{lstlisting}

\subsubsection{Run Application}
\begin{lstlisting}[language=bash]
python model_comparison_app.py
\end{lstlisting}

\subsubsection{Access}
Open browser to: \texttt{http://localhost:7860}

\section{Configuration}

\subsection{CPU vs GPU}

The application is configured to use \textbf{CPU by default} for deployment compatibility:

\begin{lstlisting}[language=Python]
# CPU inference (default)
results = model(image, device='cpu')
\end{lstlisting}

To enable GPU locally (if available):
\begin{lstlisting}[language=Python]
# GPU inference
results = model(image, device='cuda')
\end{lstlisting}

\subsection{Model Paths}

Update model paths in \texttt{model\_comparison\_app.py}:
\begin{lstlisting}[language=Python]
model1_path = "weights/best.pt"
model2_path = "VisionSentry-dex-rgb-model-replication/weights/best.pt"
\end{lstlisting}

\subsection{Port Configuration}

Default port is 7860. To change:
\begin{lstlisting}[language=Python]
demo.launch(server_port=8080)
\end{lstlisting}

\section{Troubleshooting}

\subsection{Common Issues}

\subsubsection{Out of Memory}
\begin{tcolorbox}[colback=yellow!5!white,colframe=yellow!75!black]
\textbf{Solution:} Reduce video resolution or process fewer frames at once.
\end{tcolorbox}

\subsubsection{Slow Inference}
\begin{tcolorbox}[colback=yellow!5!white,colframe=yellow!75!black]
\textbf{Solution:} CPU inference is slower. Consider upgrading to a paid tier with GPU support on Hugging Face Spaces.
\end{tcolorbox}

\subsubsection{Model Not Found}
\begin{tcolorbox}[colback=yellow!5!white,colframe=yellow!75!black]
\textbf{Solution:} Ensure model weights are uploaded and paths are correct.
\end{tcolorbox}

\subsection{Platform-Specific Issues}

\subsubsection{Hugging Face Spaces}
\begin{itemize}
    \item \textbf{Build fails:} Check \texttt{requirements\_app.txt} for incompatible versions
    \item \textbf{App crashes:} Check logs in Space settings
    \item \textbf{Slow loading:} Upgrade to better hardware tier
\end{itemize}

\subsubsection{Railway/Render}
\begin{itemize}
    \item \textbf{Deployment timeout:} Large model files may cause timeout. Use Git LFS
    \item \textbf{Port binding:} Ensure app binds to \texttt{0.0.0.0}, not \texttt{localhost}
\end{itemize}

\section{Performance Optimization}

\subsection{Reduce Model Size}
\begin{itemize}
    \item Use smaller YOLO models (YOLOv11n instead of YOLOv11x)
    \item Quantize models for faster inference
    \item Use ONNX format for optimized inference
\end{itemize}

\subsection{Video Processing}
\begin{itemize}
    \item Limit video length (e.g., max 30 seconds)
    \item Reduce frame rate (process every Nth frame)
    \item Lower output resolution
\end{itemize}

\subsection{Caching}
\begin{itemize}
    \item Cache model loading
    \item Use Gradio's built-in caching
    \item Implement result caching for repeated inputs
\end{itemize}

\section{Security Considerations}

\subsection{Best Practices}
\begin{enumerate}
    \item \textbf{Input validation:} Limit file sizes and formats
    \item \textbf{Rate limiting:} Prevent abuse with request limits
    \item \textbf{HTTPS:} All platforms provide automatic HTTPS
    \item \textbf{Environment variables:} Store sensitive data securely
\end{enumerate}

\subsection{File Upload Limits}
\begin{lstlisting}[language=Python]
# Limit video size to 100MB
video_input = gr.Video(label="Upload Video", 
                       file_count="single",
                       max_size=100*1024*1024)
\end{lstlisting}

\section{Monitoring and Maintenance}

\subsection{Logging}
Enable logging for debugging:
\begin{lstlisting}[language=Python]
import logging
logging.basicConfig(level=logging.INFO)
\end{lstlisting}

\subsection{Usage Analytics}
\begin{itemize}
    \item Hugging Face Spaces provides built-in analytics
    \item Railway/Render offer usage dashboards
    \item Monitor CPU/memory usage
\end{itemize}

\subsection{Updates}
\begin{enumerate}
    \item Test changes locally first
    \item Use version control (Git)
    \item Deploy to staging before production
    \item Monitor after deployment
\end{enumerate}

\section{Cost Estimation}

\subsection{Free Tiers}
\begin{itemize}
    \item \textbf{Hugging Face Spaces:} Free CPU tier (2 vCPU, 16GB RAM)
    \item \textbf{Railway:} \$5 free credit/month
    \item \textbf{Render:} 750 hours/month free
\end{itemize}

\subsection{Paid Tiers (if needed)}
\begin{itemize}
    \item \textbf{Hugging Face GPU:} \$0.60/hour (T4 GPU)
    \item \textbf{Railway:} \$0.000231/GB-second
    \item \textbf{Render:} \$7/month (starter)
\end{itemize}

\section{Quick Start Checklist}

\begin{enumerate}
    \item[$\square$] Choose deployment platform (Hugging Face recommended)
    \item[$\square$] Create account on chosen platform
    \item[$\square$] Prepare files: \texttt{app.py}, \texttt{requirements\_app.txt}, model weights
    \item[$\square$] Upload files to platform
    \item[$\square$] Configure settings (CPU, port, etc.)
    \item[$\square$] Deploy and wait for build
    \item[$\square$] Test application with sample video
    \item[$\square$] Share public URL
\end{enumerate}

\section{Support and Resources}

\subsection{Documentation}
\begin{itemize}
    \item Gradio: \url{https://gradio.app/docs}
    \item Hugging Face Spaces: \url{https://huggingface.co/docs/hub/spaces}
    \item Ultralytics YOLO: \url{https://docs.ultralytics.com}
\end{itemize}

\subsection{Community}
\begin{itemize}
    \item Gradio Discord: \url{https://discord.gg/gradio}
    \item Hugging Face Forums: \url{https://discuss.huggingface.co}
\end{itemize}

\section{Conclusion}

This guide covered multiple deployment options for the Drone Detection Model Comparison app. For most users, \textbf{Hugging Face Spaces} is the recommended platform due to its ease of use, ML optimization, and generous free tier.

For production deployments with high traffic, consider upgrading to paid tiers or using cloud platforms like Google Cloud Run or AWS with auto-scaling capabilities.

\vspace{1cm}
\noindent\textbf{Happy Deploying!}

\end{document}
